\documentclass[11pt,a4paper]{article}
\usepackage[T1]{fontenc}
\usepackage[utf8]{inputenc}
\usepackage{lmodern,amsmath,amssymb}
\usepackage[margin=25mm]{geometry}
\usepackage[hidelinks]{hyperref}
\hypersetup{pdftitle={Openings for the mass gap, and why it is not special to SU(3)},pdfauthor={navstokgap research working notes}}
\newcommand{\tightlist}{\setlength{\itemsep}{0pt}\setlength{\parskip}{0pt}}
\setlength{\emergencystretch}{3em}
\setcounter{secnumdepth}{0}
\begin{document}
\hypertarget{openings-for-the-mass-gap-and-why-it-is-not-special-to-su3}{%
\section{Openings for the mass gap, and why it is not special to
SU(3)}\label{openings-for-the-mass-gap-and-why-it-is-not-special-to-su3}}

Recorded 2026-09-23 at the user's request and \textbf{parked}: the
Planck-gap paper comes first, and none of these is being attacked now,
except possibly the fourth, which is the link from the Planck-gap work.
The note answers two questions the user asked and corrects one statement
of the map.

\textbf{The gap is expected for every compact simple non-abelian group.}
Lattice evidence shows gapped glueball spectra for \(SU(N)\) from
\(N=2\) to large \(N\), with masses smooth in \(1/N^2\), and for the
other classical groups; the dividing line is between abelian and
non-abelian, since compact \(U(1)\) in four dimensions has a proven
massless phase at weak coupling. \(SU(3)\) fixes the constants; a method
that works will almost surely be group-general.

\textbf{A precision on the map.} The
\href{mass-gap-position.md}{position note} says that the conjecture is
T2\('\) together with T3, given T4, and the
\href{mass-gap-obligations-lattice.md}{obligations map} says that T3
presupposes T2\('\). T3 needs less: a gap \(\delta_\infty(g)>0\) on some
interval \((0,g_1)\) along the scaling curve. T2\('\), a gap at every
lattice coupling, is a property of a particular lattice action, and for
the Wilson action it fails for \(SU(N\ge5)\), where a first-order bulk
transition occurs at intermediate coupling. For \(SU(3)\) with the
Wilson action the bulk behaviour is a crossover, so T2\('\) remains
plausible there; the conditional theorem of the
\href{mass-gap-conditional-theorem.md}{conditional-theorem note}, stated
along the trajectory, needs no T2\('\) at all.

\textbf{Four openings,} ranked by how directly they reach the place
where the map ends, the control of about six blocking steps at order-one
coupling.

\begin{enumerate}
\def\labelenumi{\arabic{enumi}.}
\tightlist
\item
  Curvature: the Bakry--Émery criterion and its multiscale form.
\item
  The centre-stabilized theory on \(\mathbb R^3\times S^1\) at small
  radius.
\item
  Three dimensions first.
\item
  The Planck-gap link: Simon's valley lifting, a floor that holds for
  every state.
\end{enumerate}

\hypertarget{evidence-across-groups}{%
\subsection{1. Evidence across groups}\label{evidence-across-groups}}

\begin{itemize}
\tightlist
\item
  \textbf{\(SU(N)\).} Glueball masses and string tensions computed for
  \(N\) from \(2\) to about \(12\) are finite and extrapolate smoothly
  in \(1/N^2\) (Lucini and Panero,
  \href{https://doi.org/10.1016/j.physrep.2013.01.001}{Phys.
  Rep.~\textbf{526}, 93, 2013}, metadata; Lucini and Teper,
  \href{https://doi.org/10.1088/1126-6708/2001/06/050}{JHEP
  \textbf{2001}(06), 050}, metadata).
\item
  \textbf{\(G_2\).} The exceptional group has a trivial centre and still
  shows confinement behaviour and a first-order deconfinement transition
  (Pepe and Wiese,
  \href{https://doi.org/10.1016/j.nuclphysb.2006.12.024}{Nucl. Phys. B
  \textbf{768}, 21, 2007}, metadata), so the centre is not what the gap
  rests on.
\item
  \textbf{\(U(1)\).} Four-dimensional compact \(U(1)\) has a massless
  Coulomb phase at weak coupling (Guth,
  \href{https://doi.org/10.1103/PhysRevD.21.2291}{PRD \textbf{21}, 2291,
  1980}; Fröhlich and Spencer,
  \href{https://doi.org/10.1007/BF01213610}{CMP \textbf{83}, 411, 1982};
  metadata), which is why the \href{abelian-misses-the-box.md}{abelian
  note} separates the groups by the running, \(b_0=0\) against
  \(b_0=11N/(48\pi^2)\). In three dimensions compact \(U(1)\) is gapped
  at every coupling (Polyakov,
  \href{https://doi.org/10.1016/0550-3213(77)90086-4}{Nucl. Phys. B
  \textbf{120}, 429, 1977}; proved by Göpfert and Mack,
  \href{https://doi.org/10.1007/BF01961240}{CMP \textbf{82}, 545, 1982};
  metadata), with a gap that vanishes in physical units in the
  continuum.
\end{itemize}

\textbf{Bulk transitions depend on the action.} With the standard Wilson
plaquette action, ``the bulk transition is a rapid cross-over if
\(N\le4\) but becomes an increasingly strong first order transition for
\(N\ge5\)''; for \(N=3\) a negative adjoint term avoids the bulk phase
altogether (Rindlisbacher, Rummukainen and Salami,
\href{https://arxiv.org/abs/2306.14319}{arXiv:2306.14319}, passage). In
the plane of fundamental and adjoint couplings, a line of first-order
transitions ends at a critical point near the Wilson axis (Bhanot and
Creutz, \href{https://doi.org/10.1103/PhysRevD.24.3212}{PRD \textbf{24},
3212, 1981}, metadata). By Corollary 3 of the
\href{gapped-set-critical-coupling.md}{gapped-set note}, a first-order
transition at \(g_b\) gives \(\delta_\infty(g_b)=0\) through vacuum
coexistence, so the gapped set \(\mathcal G\) has a hole there while the
physical gap, which lives at \(g\to0\), is untouched. The statement the
conjecture needs is \(\mathcal G\supseteq(0,g_1)\) for some \(g_1>0\),
joined to T3 along the scaling curve.

\hypertarget{why-nge5-differs-and-what-sits-beside-the-su3-wilson-axis}{%
\subsubsection{\texorpdfstring{Why \(N\ge5\) differs, and what sits
beside the \(SU(3)\) Wilson
axis}{Why N\textbackslash ge5 differs, and what sits beside the SU(3) Wilson axis}}\label{why-nge5-differs-and-what-sits-beside-the-su3-wilson-axis}}

Give the lattice action a fundamental coupling \(\beta_F\) and an
adjoint coupling \(\beta_A\). In the \((\beta_F,\beta_A)\) plane a line
of first-order bulk transitions ends at a critical point (Bhanot and
Creutz 1981). For small \(N\) the Wilson axis \(\beta_A=0\) passes
beside the endpoint and shows a rapid crossover; from \(N=5\) it crosses
the line, and the transition strengthens with \(N\). The large-\(N\)
limit makes a non-analyticity between strong and weak lattice coupling
unavoidable. In two dimensions, where the lattice theory is exactly
solvable, it is analytic at every finite \(N\) and has a third-order
transition at a critical 't Hooft coupling at \(N=\infty\) (Gross and
Witten, \href{https://doi.org/10.1103/PhysRevD.21.446}{PRD \textbf{21},
446, 1980}; Wadia,
\href{https://doi.org/10.1016/0370-2693(80)90353-6}{PLB \textbf{93},
403, 1980}; metadata). In four dimensions the large-\(N\) transition is
first order, and at finite \(N\) it survives as a genuine transition
from \(N=5\) and is smoothed into a crossover below.

For \(SU(3)\) the endpoint matters twice.

\begin{itemize}
\tightlist
\item
  \textbf{It realizes the second-order branch.} Approaching the
  endpoint, ``the mass of the \(0^{++}\) glueball decreases \ldots{}
  while the other glueball masses appear unchanged'', which is
  ``consistent with the notion that the bulk transition line ends in a
  critical endpoint with the continuum limit there being a \(\phi^4\)
  theory with a diverging correlation length only in the \(0^{++}\)
  channel'' (Heller,
  \href{https://doi.org/10.1016/0370-2693(95)01186-T}{PLB \textbf{362},
  123, 1995};
  \href{https://arxiv.org/abs/hep-lat/9508009}{hep-lat/9508009},
  abstract). So lattice \(SU(3)\) has a second continuum limit at a
  finite point of its two-coupling plane, a scalar theory, and the
  statement that the asymptotically free limit is the only one concerns
  families that avoid that point, such as the Wilson axis.
\item
  \textbf{It sits beside the point where the map places H2.} The
  \href{confinement-scale-bands.md}{bands note} scaled the continuum
  ratio \(r_0m_{0^{++}}=4.21\) to \(\beta_W=5.7\) and found
  \(\xi/a=0.70\). Near there the Wilson axis passes beside the endpoint,
  where the scalar glueball softens, so the lattice \(\xi/a\) at
  \(\beta_W\simeq5.7\) is likely larger, and the measured lattice value
  should replace the scaled one. The conditional theorem allows the
  remedy: blocking generates adjoint and higher-representation terms
  anyway, so the reference interaction \(\Phi_0\) of H2 can be taken at
  a small negative adjoint coupling, away from the endpoint. For \(N=3\)
  such an action avoids the bulk phase, as Rindlisbacher, Rummukainen
  and Salami report from earlier work (arXiv:2306.14319, passage).
\end{itemize}

\hypertarget{what-the-problem-statement-itself-points-to}{%
\subsubsection{What the problem statement itself points
to}\label{what-the-problem-statement-itself-points-to}}

Jaffe and Witten's official description, held as a
\href{../docs/JaffeWitten_YangMills.pdf}{local PDF}, names in its §6.6
four routes to the gap (passage), and two of them are openings 1 and 4
below:

\begin{itemize}
\tightlist
\item
  a ``duality transformation'', with Debye screening in the Coulomb gas,
  proved through the sine-Gordon representation, as the model;
\item
  the quartic potential \((A\wedge A)^2\) (Feynman,
  \href{https://doi.org/10.1016/0550-3213(81)90005-5}{Nucl. Phys. B
  \textbf{188}, 479, 1981}), which ``may be tied to curvature in the
  space of connections'' (Singer,
  \href{https://doi.org/10.1088/0031-8949/24/5/002}{Phys. Scripta
  \textbf{24}, 817, 1981}): opening 1;
\item
  ``certain quantum mechanics problems with potentials involving flat
  directions \ldots{} do lead to bound states'' (Simon 1983): opening 4;
\item
  the \(1/N\) expansion ('t Hooft 1974): gauge theory with \(SU(N)\),
  \(SO(N)\) or \(Sp(N)\) ``may be equivalent to a string theory with
  \(1/N\) as the string coupling constant'', which ``might give a
  clear-cut explanation of the mass gap and confinement and perhaps a
  good starting point for a rigorous proof (for sufficiently large
  \(N\))''. Maldacena's paper is cited there as ``surprising progress
  along these lines for certain strongly coupled four-dimensional gauge
  systems with matter, but as of yet there is no effective approach to
  the gauge theory without fermions.'' Earlier, §6.1 allows that ``there
  might some day be an asymptotic solution in a large \(N\) limit.''
\end{itemize}

\textbf{Chronology.} Maldacena's conjecture was posted in November 1997
(\href{https://arxiv.org/abs/hep-th/9711200}{hep-th/9711200}); Witten's
holographic account of non-supersymmetric Yang--Mills, in which ``the
spontaneous breaking of the center of the gauge group, magnetic
confinement, and the mass gap are coded in classical geometry'', in
March 1998 (\href{https://arxiv.org/abs/hep-th/9803131}{hep-th/9803131},
abstract); supergravity glueball masses in agreement with lattice ratios
in June 1998 (Csáki, Ooguri, Oz and Terning,
\href{https://arxiv.org/abs/hep-th/9806021}{hep-th/9806021}, abstract);
the Clay problems in May 2000. The statement was written with holography
in hand, by one of its authors, and records no effective approach to the
pure gauge theory. The usual reasons, not checked here against a primary
text, are that the tractable supergravity regime is the opposite limit
from the continuum one, with Kaluza--Klein modes at the glueball scale,
and that it holds at large \(N\) only.

\textbf{The large-\(N\) thread.} It runs through Witten's own work: the
\(1/N\) expansion for baryons
(\href{https://doi.org/10.1016/0550-3213(79)90232-3}{Nucl. Phys. B
\textbf{160}, 57, 1979}) and for the \(U(1)\) problem
(\href{https://doi.org/10.1016/0550-3213(79)90031-2}{Nucl. Phys. B
\textbf{156}, 269, 1979}), and the large-\(N\) lattice transition with
Gross (1980), the non-analyticity behind the \(N\ge5\) bulk transition
above. This programme records the route and does not adopt it: the
conjecture is posed for every compact simple group, so a proof for
sufficiently large \(N\) would leave \(SU(3)\) open unless it came with
explicit control down to \(N=3\); and at large \(N\) a lattice route
meets the bulk transition, which it would have to avoid with a
bulk-preventing action or with the centre-stabilized volume independence
of opening 2, itself a large-\(N\) construction.

\textbf{The panorama with \(\hbar\) explicit.} This repository writes
the action as
\(S=-\frac{\hbar}{4g^2}\int F^a_{\mu\nu}F^{a\mu\nu}\,d^3x\,c\,dt\) with
\(g\) dimensionless (\href{low-dimensional-mass-gap.md}{G07},
Proposition 7), so the weight \(e^{iS/\hbar}\) contains \(g\) alone and
\(\hbar\) is hidden inside it: with a classical coupling \(g_{\rm cl}\),
\(S=-\frac1{4g_{\rm cl}^2}\int\cdots\), one has
\(g^2=\hbar g_{\rm cl}^2\). Three consequences follow.

\begin{itemize}
\tightlist
\item
  \emph{The 't Hooft limit is a classical limit.} A diagram with \(E\)
  propagators, \(V\) vertices and genus \(h\) carries
  \((\hbar\lambda_{\rm cl})^{E-V}N^{2-2h}\) with
  \(\lambda_{\rm cl}=g_{\rm cl}^2N\) ('t Hooft,
  \href{https://doi.org/10.1016/0550-3213(74)90154-0}{Nucl. Phys. B
  \textbf{72}, 461, 1974}, metadata). The loop parameter is
  \(\lambda=\hbar g_{\rm cl}^2N\), so \(N\to\infty\) at fixed
  \(\lambda\) is \(\hbar\to0\) with \(\hbar N\) fixed at fixed
  \(g_{\rm cl}\), and \(1/N^2=(\hbar g_{\rm cl}^2/\lambda)^2\): the
  genus expansion is an expansion in \(\hbar^2\). Yaffe formulates
  large-\(N\) limits in general as classical mechanics
  (\href{https://doi.org/10.1103/RevModPhys.54.407}{Rev.~Mod. Phys.
  \textbf{54}, 407, 1982}, metadata). In holography it is literal: the
  bulk Newton constant in anti-de Sitter units is of order \(1/N^2\), so
  the bulk loop expansion is the \(1/N^2\) expansion.
\item
  \emph{The gap is non-perturbative in \(\hbar\).} With the two-loop
  scale of the \href{mass-gap-obligations-lattice.md}{obligations map}
  §5 and \(b_0=11N/(48\pi^2)\),
  \(mc^2\propto\frac{\hbar c}a\exp\bigl(-\frac{24\pi^2}{11\,\hbar g_{\rm cl}^2N}\bigr)\)
  times a power of \(\hbar g_{\rm cl}^2N\): at fixed \(N\), \(a\) and
  \(g_{\rm cl}\) it vanishes faster than any power of \(\hbar\), the way
  a tunnelling splitting does, which gives one precise reading of Jaffe
  and Witten's remark that the gap is ``not classically visible''. It
  depends on \(\hbar\) and \(N\) only through \(\hbar N\), so the 't
  Hooft limit is the classical limit along which the gap survives.
\item
  \emph{The small volume is a power of \(\hbar\).} G07's zero-mode gap
  \(\delta_1g^{2/3}\hbar c/L\) is
  \(\delta_1\hbar^{4/3}g_{\rm cl}^{2/3}c/L\): a floor of the Planck-gap
  kind, a power of \(\hbar\) from the zero-point energy that lifts the
  valleys. The two regimes meet at the crossover \(z\simeq2\) of the
  \href{torus-valley-potential.md}{valley note}. Opening 4 lives in the
  power-law regime, where the Planck-gap methods apply; the
  infinite-volume gap lives in the non-perturbative one.
\end{itemize}

\hypertarget{the-openings}{%
\subsection{2. The openings}\label{the-openings}}

\hypertarget{opening-1-curvature}{%
\subsubsection{Opening 1: curvature}\label{opening-1-curvature}}

The configuration space of lattice gauge theory is a product of copies
of the group, and a compact simple group carries positive Ricci
curvature: for \(SU(N)\) with the metric
\(\langle X,Y\rangle=\operatorname{Tr}(XY^*)\),
\(\operatorname{Ric}=\frac N2\,g\). The Bakry--Émery criterion turns a
positive lower bound on \(\operatorname{Ric}+\operatorname{Hess}S\) into
a log-Sobolev and a Poincaré inequality, that is into a spectral gap of
the Langevin dynamics and exponential decay of correlations (Bakry and
Émery, \href{https://doi.org/10.1007/BFb0075847}{LNM 1123, 177, 1985},
metadata). A flat group, \(U(1)\), gets nothing from the first term,
which lines up with the abelian failure.

Shen, Zhu and Zhu use exactly this to prove, for lattice Yang--Mills
with 't Hooft coupling \(\beta N\) in any dimension \(d>1\), uniqueness
of the infinite-volume measure, log-Sobolev and Poincaré inequalities
and exponential decay of correlations, ``a strictly positive mass gap'',
for \(|\beta|<1/(16(d-1))\) when the group is \(SU(N)\)
(\href{https://doi.org/10.1007/s00220-022-04609-1}{CMP \textbf{400},
2022}; \href{https://arxiv.org/abs/2204.12737}{arXiv:2204.12737},
abstract). In the Wilson normalization of this repository,
\(\beta_W=2N/g^2\) with weight
\(e^{(\beta_W/N)\sum\operatorname{Re}\operatorname{Tr}U_p}\), their
\(\beta\) is \(\beta_W/N^2\), so for \(SU(3)\) in four dimensions the
condition is \(\beta_W<3/16\), that is

\[g^2>32 .\]

That conversion is ours and must be checked against their definitions.
If it stands, the curvature method already beats this repository's
strong-coupling thresholds, \(g^2\ge176\) for the Wilson transfer matrix
(\href{wilson-strong-coupling-explicit.md}{note}), \(g^2\ge388\) for
Kogut--Susskind
(\href{kogut-susskind-strong-coupling-explicit.md}{note}) and
\(g^2>444\) by Dobrushin (\href{dobrushin-uniqueness-wilson.md}{note}),
and narrows band B of the \href{confinement-scale-bands.md}{bands note}
from \(\beta_W\in(0.0135,5.7)\) to \(\beta_W\in(0.19,5.7)\).

The weak side is where the multiscale form enters. Bauerschmidt and
Bodineau derive ``a multiscale generalisation of the Bakry--Émery
criterion'' through the Polchinski renormalization equation, ``effective
for measures which are far from log-concave'', and prove optimal
log-Sobolev inequalities for the continuum sine-Gordon model at
\(\beta<6\pi\) (\href{https://doi.org/10.1002/cpa.21926}{CPAM
\textbf{74}, 2020};
\href{https://arxiv.org/abs/1907.12308}{arXiv:1907.12308}, abstract).
Its output is a spectral gap controlled at every scale, which is the
form of hypothesis H1 in the conditional theorem. It also sits outside
the family closed in
\href{flow-before-decimation.md}{flow-before-decimation}: that note
closes every criterion that depends on the perturbation through a norm,
and a curvature bound is a lower bound on a Hessian.

The risk is concrete. After a blocking step the effective action is a
function on the coarse configuration manifold, and the criterion needs
\(\operatorname{Ric}+\operatorname{Hess}S_{\rm eff}\) bounded below at
every scale. At weak coupling the Hessian is degenerate along gauge
orbits and along the abelian valleys, and the curvature of the group has
to carry those directions. \textbf{First note if opened:} check the
conversion to \(g^2>32\); compute the curvature budget for \(SU(2)\) and
\(SU(3)\) against the Hessian of the Wilson action; and state exactly
what lower bound one blocking step must preserve.

\hypertarget{opening-2-the-centre-stabilized-theory-at-small-radius}{%
\subsubsection{Opening 2: the centre-stabilized theory at small
radius}\label{opening-2-the-centre-stabilized-theory-at-small-radius}}

Ünsal and Yaffe add a double-trace deformation that keeps centre
symmetry unbroken; on \(\mathbb R^3\times S^1\) at small radius ``the
deformed Yang--Mills theory has a mass gap and exhibits linear
confinement'', by an analytic semiclassical treatment, and ``there are
no order parameters which distinguish the small and large radius
regimes'' (\href{https://doi.org/10.1103/PhysRevD.78.065035}{PRD
\textbf{78}, 065035, 2008}, abstract). The mechanism at small radius is
Polyakov's monopole plasma, which Göpfert and Mack made rigorous for
three-dimensional compact \(U(1)\) at every coupling. \textbf{Target if
opened:} a rigorous mass gap for the deformed theory at small radius,
with explicit constants. That would be a theorem about a deformed
theory; carrying it to \(\mathbb R^4\) is the adiabatic-continuity
conjecture.

\hypertarget{opening-3-three-dimensions-first}{%
\subsubsection{Opening 3: three dimensions
first}\label{opening-3-three-dimensions-first}}

In three dimensions \(g^2\) has the dimensions of energy, the theory is
superrenormalizable, and the coupling at lattice scale is
\(g^2a/\hbar c\), which vanishes as a power of \(a\). The weak-side
deficit of twelve orders of magnitude found in
\href{small-field-step-decay-and-threshold.md}{the small-field threshold
note} comes from the logarithmic running of four dimensions, and in
three dimensions the ultraviolet steps shrink geometrically, so the
constant-chasing that failed in four dimensions may close. The
\href{three-dimensional-gap-one-function.md}{one-function note} reduces
the problem to a lower bound \(f(x)\ge f_->0\) on one function of one
variable. A three-dimensional gap uniform in the spacing is open and
would be a major result by itself.

\hypertarget{opening-4-the-planck-gap-link}{%
\subsubsection{Opening 4: the Planck-gap
link}\label{opening-4-the-planck-gap-link}}

The small-volume mechanism of \href{low-dimensional-mass-gap.md}{G07}
and \href{action-floor-yang-mills-gap.md}{G08} is Simon's: along the
flat valleys of the commutator potential the transverse modes are
oscillators whose frequency grows with the distance along the valley,
and their zero-point energy lifts the valley (Simon,
\href{https://doi.org/10.1016/0003-4916(83)90057-X}{Ann. Phys.
\textbf{146}, 209, 1983}, metadata). In its simplest form, for fixed
\(x\),

\[\frac{\hat p_y^2}{2m}+\frac12m\Omega(x)^2\hat y^2\ \ge\ \frac12\hbar\,\Omega(x),\]

an operator inequality that holds in every state. The
\href{torus-valley-potential.md}{torus-valley note} finds the
field-theory version term by term: the linear term \(2|a|\) of the
valley potential is the zero-point energy of the \(k=0\) charged modes.

The Planck-gap work sharpens what kind of floor this is. Its floors that
rest on second moments of a prepared state fall to squeezing and to grid
states, and the ones that survive rest on the commutator alone: the
phase between Newton's polygon and the parabola, and the cost of
disturbance. Simon's inequality is of the surviving kind, a spectral
statement, so the valley lifting is the mass-gap instance of the lesson
the Planck paper draws. \textbf{First step if opened:} hypothesis (H3)
of the \href{weak-coupling-feshbach-reduction.md}{Feshbach note}, the
gap of the zero-mode Hamiltonian \(h_3\) plus the even valley potential
\(U(a)\), proved by a Simon-type operator inequality with every constant
explicit. That is the lower side of the small-volume corner S at weak
coupling, and it stays within reach of the Planck-gap methods.

\hypertarget{consequence-for-state}{%
\subsection{3. Consequence for STATE}\label{consequence-for-state}}

The mass-gap track stays paused. This note records the four openings and
the answer on groups: the gap is expected for every compact simple
non-abelian group, and \(SU(3)\) is the group that fixes the constants.
The map's statement that the conjecture is T2\('\) with T3 is too
strong: T3, with a gap on an interval \((0,g_1)\) along the scaling
curve, is what is needed, and T2\('\) fails for \(SU(N\ge5)\) with the
Wilson action. If one opening is taken up next, it is the fourth, which
uses the Planck-gap methods directly; the first carries the most
promise, and its first check is the conversion of Shen, Zhu and Zhu's
threshold to \(g^2>32\).
\end{document}
